\chapter{Testovanie a vyhodnotenie}
\label{ch:testovanie}

Systém bol overený na troch úrovniach: (1)~automatické testy softvéru, (2)~meranie
presnosti jednotlivých biometrických modulov podľa noriem a (3)~end-to-end test celého
systému s~reálnymi používateľmi a útokmi. Výkon bol meraný samostatne. Všetky
vyhodnocovacie skripty sú v~priečinku \texttt{backend/evaluation} a výsledky ukladajú
ako tabuľky (JSON, CSV, Markdown) a grafy (PNG, PDF).

\section{Metriky}

\subsection*{Rozpoznávanie tváre}
Pri prahu $\tau$ a skóre porovnaní genuine $G$ (rovnaká osoba) a impostor $I$ (rôzne
osoby):
\begin{equation}
  \mathrm{FAR}(\tau) = \frac{|\{s \in I : s \ge \tau\}|}{|I|},
  \qquad
  \mathrm{FRR}(\tau) = \frac{|\{s \in G : s < \tau\}|}{|G|}.
\end{equation}
Miera rovnakej chyby (EER) je hodnota v~bode $\mathrm{FAR}(\tau) = \mathrm{FRR}(\tau)$;
počíta sa lineárnou interpoláciou medzi susednými prahmi. Ďalej sa uvádza FRR pri
pevnom FAR (1~\% a 0,1~\%) a podiel snímok bez detegovanej tváre (FTA).

\subsection*{Detekcia prezentačných útokov}
Podľa ISO/IEC 30107-3 \cite{iso30107-3, busch2023standards} sa pre každý druh útoku
(\emph{PAI species}) $k$ počíta
\begin{equation}
  \mathrm{APCER}_k = \frac{\text{počet útokov typu } k \text{ prijatých ako živá tvár}}{\text{počet útokov typu } k},
  \qquad
  \mathrm{BPCER} = \frac{\text{počet živých tvárí odmietnutých}}{\text{počet živých tvárí}},
\end{equation}
pričom celkové $\mathrm{APCER} = \max_k \mathrm{APCER}_k$ (najhorší prípad) a
$\mathrm{ACER} = (\mathrm{APCER} + \mathrm{BPCER})/2$. Uvádza sa aj BPCER20 a BPCER100,
teda BPCER pri prahu, pri ktorom $\mathrm{APCER} \le 5~\%$ resp. $1~\%$.

\section{Automatické testy}

Serverová časť má 47 testov (pytest) s~pokrytím kódu 87~\%. Testy nahrádzajú neurónové
siete deterministickými náhradami: syntetická snímka nesie v~jednom pixeli identitu,
polohu hlavy a príznak živosti. Vďaka tomu sa celý tok overuje cez skutočné HTTP
rozhranie bez modelov a bez GPU. Klient má jednotkové testy v~nástroji Vitest. Testy,
statická analýza (Ruff, mypy, ESLint) a zostavenie Docker obrazov bežia automaticky
pri každej zmene v~GitHub Actions.

\begin{table}[H]
  \centering
  \caption{Pokrytie automatickými testami}
  \label{tab:testy}
  \small
  \begin{tabularx}{\textwidth}{@{}l X@{}}
    \toprule
    Oblasť & Overované správanie \\
    \midrule
    Póza hlavy & znamienko yaw a pitch, nezávislosť od náklonu hlavy, degenerované body \\
    Aktívna výzva & správny a nesprávny smer, nedostatočný pohyb, málo snímok s~tvárou, náhodnosť poradia \\
    Metriky & FAR/FRR, EER (aj analyticky pre dve normálne rozdelenia), APCER/BPCER/ACER \\
    Bezpečnosť & podpis a platnosť JWT, šifrovanie vzorov, produkčná konfigurácia, rate-limit \\
    Výzvy & jednorazovosť, expirácia, kontrola účelu \\
    End-to-end API & registrácia a overenie, impostor, neznámy používateľ, útok fotografiou (pasívny aj aktívny), zámena osoby počas relácie, opakované použitie výzvy, chýbajúci krok, duplicitná tvár, identifikácia 1:N, zablokovanie účtu, výmaz účtu, práva administrátora \\
    \bottomrule
  \end{tabularx}
\end{table}

\section{Presnosť rozpoznávania tváre}

\textbf{Postup.} Model SFace bol vyhodnotený na databáze Labeled Faces in the Wild
podľa protokolu \emph{pairs} (6000 dvojíc: 3000 genuine, 3000 impostor) príkazom
\texttt{python -m evaluation.evaluate\_recognition --lfw}. Kontrola kvality bola pri
tomto meraní vypnutá, aby sa meral samotný model rozpoznávania.

\todo{Spustite skript a doplňte hodnoty z~\texttt{backend/reports/recognition-*/table.md};
grafy \texttt{far\_frr.pdf}, \texttt{det.pdf}, \texttt{score\_distribution.pdf} skopírujte
do \texttt{report/figures/}.}

\begin{table}[H]
  \centering
  \caption{Výsledky rozpoznávania tváre na LFW}
  \label{tab:lfw}
  \begin{tabular}{@{}lccc@{}}
    \toprule
    Pracovný bod & Prah & FAR [\%] & FRR [\%] \\
    \midrule
    EER & \todo{} & \todo{} & \todo{} \\
    Systémový prah & 0,40 & \todo{} & \todo{} \\
    FAR $\le$ 1~\% & \todo{} & \todo{} & \todo{} \\
    FAR $\le$ 0,1~\% & \todo{} & \todo{} & \todo{} \\
    \bottomrule
  \end{tabular}
\end{table}

\begin{figure}[H]
  \centering
  \begin{subfigure}{0.49\textwidth}
    \centering
    \optfigure{score_distribution.pdf}{\linewidth}{histogram skóre}
    \caption{Rozdelenie skóre}
  \end{subfigure}
  \hfill
  \begin{subfigure}{0.49\textwidth}
    \centering
    \optfigure{det.pdf}{\linewidth}{DET krivka}
    \caption{DET krivka}
  \end{subfigure}
  \caption{Rozpoznávanie tváre na LFW}
  \label{fig:lfw}
\end{figure}

\todo{Interpretácia (vlastnými slovami): ako ďaleko sú rozdelenia genuine a impostor,
čo znamená systémový prah pre bezpečnosť a pohodlie, porovnanie s~hodnotou EER.}

\section{Presnosť detekcie prezentačných útokov}

\textbf{Dataset.} Pomocou skriptu \texttt{evaluation.capture\_dataset} bol webkamerou
nahraný vlastný dataset s~nasledujúcimi triedami:

\begin{table}[H]
  \centering
  \caption{Dataset pre vyhodnotenie PAD}
  \label{tab:pad-dataset}
  \begin{tabularx}{\textwidth}{@{}l X r@{}}
    \toprule
    Trieda & Popis & Snímky \\
    \midrule
    \texttt{bona\_fide} & živé tváre autorov, rôzne vzdialenosti, svetlo a uhly & \todo{} \\
    \texttt{attack/print} & fotografia vytlačená na papieri A4 & \todo{} \\
    \texttt{attack/replay} & fotografia a video na displeji mobilného telefónu & \todo{} \\
    \texttt{attack/screen} & fotografia na monitore notebooku & \todo{} \\
    \bottomrule
  \end{tabularx}
\end{table}

\textbf{Postup.} Príkaz \texttt{python -m evaluation.evaluate\_pad --dir data/pad}
vypočíta skóre živosti každej snímky a z~nich metriky pri systémovom prahu 0,70.

\begin{table}[H]
  \centering
  \caption{Výsledky pasívnej detekcie útokov (prah 0,70)}
  \label{tab:pad}
  \begin{tabular}{@{}lc@{}}
    \toprule
    Metrika & Hodnota [\%] \\
    \midrule
    APCER (print) & \todo{} \\
    APCER (replay) & \todo{} \\
    APCER (screen) & \todo{} \\
    APCER (max) & \todo{} \\
    BPCER & \todo{} \\
    ACER & \todo{} \\
    EER & \todo{} \\
    BPCER20 & \todo{} \\
    \bottomrule
  \end{tabular}
\end{table}

\begin{figure}[H]
  \centering
  \begin{subfigure}{0.49\textwidth}
    \centering
    \optfigure{pad_score_distribution.pdf}{\linewidth}{skóre živosti podľa triedy}
    \caption{Rozdelenie skóre živosti}
  \end{subfigure}
  \hfill
  \begin{subfigure}{0.49\textwidth}
    \centering
    \optfigure{pad_det.pdf}{\linewidth}{DET pre každý typ útoku}
    \caption{DET (APCER vs. BPCER)}
  \end{subfigure}
  \caption{Pasívna detekcia prezentačných útokov}
  \label{fig:pad}
\end{figure}

\todo{Interpretácia: ktorý typ útoku je najťažší, či je prah 0,70 vhodný, prečo samotná
pasívna detekcia nestačí (generalizácia na nové podmienky \cite{sun2023rethinking}).}

Už pri overovaní implementácie sa ukázal dôležitý jav: digitálna fotografia odoslaná
priamo serveru dostala od MiniFASNet skóre živosti 0,999 -- pasívny model ju neodlíšil
od živej tváre, pretože na nej nie sú stopy tlače ani displeja. Takú reláciu však
zamietla aktívna výzva (všetky kroky s~nulovým pohybom). Tento prípad ilustruje, prečo
systém kombinuje pasívnu a aktívnu detekciu.

\section{End-to-end test systému}

Test prebehol na verejnej inštancii aplikácie. Výsledky sa čerpajú z~auditného záznamu
skriptom \texttt{python -m evaluation.audit\_report}.

\begin{table}[H]
  \centering
  \caption{Scenáre end-to-end testu}
  \label{tab:e2e}
  \small
  \begin{tabularx}{\textwidth}{@{}X c l c@{}}
    \toprule
    Scenár & Pokusy & Očakávaný výsledok & Úspešnosť [\%] \\
    \midrule
    Registrovaný používateľ, 1:1 & \todo{} & prijatie & \todo{} \\
    Registrovaný používateľ, 1:N & \todo{} & prijatie & \todo{} \\
    Iná osoba s~menom obete & \todo{} & \texttt{no\_match} & \todo{} \\
    Vytlačená fotografia obete & \todo{} & zamietnutie & \todo{} \\
    Fotografia na mobile/monitore & \todo{} & zamietnutie & \todo{} \\
    Video obete s~pohybmi hlavy & \todo{} & zamietnutie & \todo{} \\
    Zámena osoby počas snímania & \todo{} & \texttt{identity\_inconsistent} & \todo{} \\
    Opakované odoslanie požiadavky & \todo{} & \texttt{invalid\_challenge} & \todo{} \\
    \bottomrule
  \end{tabularx}
\end{table}

\begin{figure}[H]
  \centering
  \optfigure{rejection_reasons.pdf}{0.7\textwidth}{graf z evaluation.audit\_report}
  \caption{Dôvody zamietnutia v~end-to-end teste}
  \label{fig:reasons}
\end{figure}

\todo{Interpretácia: ktorá vrstva zachytila ktorý útok, koľko falošných zamietnutí
živých používateľov a prečo (svetlo, nedostatočný pohyb), prípadná úprava prahov.}

\section{Výkon}

Latencia bola meraná skriptom \texttt{evaluation.benchmark} na notebooku (CPU, ONNX
Runtime bez GPU) na obrázku $512\times512$~px, 30 opakovaní.

\begin{table}[H]
  \centering
  \caption{Latencia spracovania jednej snímky}
  \label{tab:latencia}
  \begin{tabular}{@{}lrrr@{}}
    \toprule
    Fáza & Priemer [ms] & Medián [ms] & P95 [ms] \\
    \midrule
    Detekcia (YuNet) & 4,97 & 4,82 & 5,28 \\
    Kontrola kvality & 0,10 & 0,10 & 0,10 \\
    Odhad polohy hlavy & 0,01 & 0,01 & 0,01 \\
    Pasívna živosť (2$\times$ MiniFASNet) & 2,44 & 2,42 & 2,56 \\
    Extrakcia príznakov (SFace) & 3,53 & 3,50 & 3,72 \\
    \midrule
    Celá snímka & 13,26 & 12,63 & 16,10 \\
    \bottomrule
  \end{tabular}
\end{table}

Relácia obsahuje 15 snímok (3 kroky po 5 snímkach), jej spracovanie na serveri teda
trvá približne 0,2~s. Celkovú dobu prihlásenia (približne 7~s) určuje najmä čas, ktorý
používateľ potrebuje na vykonanie pohybov. Na bezplatnej inštancii služby Render je
spracovanie pomalšie a prvá požiadavka po nečinnosti trvá až desiatky sekúnd, kým sa
inštancia spustí.
